\subsection[The residual pairs (1,b), b >= 6]{The residual pairs $(1,b)$, $b\ge6$}\label{loc:S7}
Put $P=qf_b$, $U_0=G_{b-1}(K)g$, $T_0=KU_0=S_b-S_1$,
$a=19/32$, $\epsilon=ad$, $\tau=\theta^4$ and
\[
 m_0=1-2ad+ad^2/2,\qquad R=T_0-\epsilon(3g+2r+s).
\]
The finite-$G$ comparison in Section~\ref{loc:S5} gives
$\|U_0\|^2\le d-2\delta_{2}/5$, $\|T_0\|^2\le \delta_{2}-2\delta_{3}/5$.
Every finite alternating truncation also satisfies
\begin{gather*}
 \E U_0\ge-d+\delta_{2}-\delta_{3},\quad
 \langle g,U_0\rangle\ge d-\delta_{2},\quad
 \langle r,U_0\rangle\ge \delta_{2}-\delta_{3},\\
 \langle g,T_0\rangle\ge-\delta_{2}+\delta_{3}-\delta_{4},\quad
 \langle r,T_0\rangle\ge \delta_{2}-\delta_{3},\quad
 \langle s,T_0\rangle\ge \delta_{3}-\delta_{4}.
\end{gather*}
These retain one negative tail term and bound the remaining decreasing
alternating sum by zero. Telescoping gives
$\E[pT_0]=-d-u_{b+1}\le-(1-\tau)d$, where $u_j=\E K^j\one$.
Retain $f_0,f_1,f_2,f_3$ in convolution, whose sum is
$4-d+3g+2r+s$. The actual envelope and mean are
\[
 q+f_b\le2m_0+R,\quad P\le(m_0+R/2)^2,\quad
 \E[pR]\le-(1-\tau)d-a^2d^2\le-31d/32.
\]
Using $R=K[U_0-\epsilon(3\one+2g+r)]$ gives
\begin{align*}
 \|R\|_\infty^2&\le d-2\delta_{2}/5+\epsilon^2(9-2d+\delta_{2})
                              +2\epsilon(d-2\delta_{2}+4\delta_{3})\le U=d+4d^2,\\
 \E R^2&\le \delta_{2}-2\delta_{3}/5+\epsilon^2(9d-2\delta_{2}+\delta_{3})
                              +2\epsilon(\delta_{2}-2\delta_{3}+4\delta_{4})\le dW,\\
 W&=11/32+33d/80+3d^2.
\end{align*}
For the first reduction the spectral part is bounded by $-2\delta_{2}/5$,
since $4-8\theta=3/4>\epsilon$; then use $\delta_{2}\ge d^2$.
For the second, the derivative of
$x-2x^2/5+\epsilon^2(9-2x+x^2)+2\epsilon(x-2x^2+4x^3)$
is positive on $[0,\theta]$; its first two parts exceed $119/200$.
Substitution at $\theta$ gives coefficients
$871/2560,53599/131072,3087633/1048576$, bounded by those in $W$.
As usual $V_p/d\le W+\sqrt{UW}$.

For $d\ge4/19$, positivity gives $f_b\le2\mu_{b-2}p$, so
$P\le2\mu_{b-2}$.
Use the path mean $\mu_n\le1-19d/8$ for $b\ge7$.
For $b=6$, the exact $\mu_4\le1-83d/32$ gives the threshold
$16/83<4/19$. Thus $P\le1$ in the entire large range.
On $d\le4/19$, comparison of $U$ with
$(1/6+4ad-ad^2)^2$ leaves at least
$1/36-5d/24+5d^2/6>0$, whose discriminant is $-85/1728$.
Therefore the global cap is $P\le(13/12)^2$, and $(13/12)^6<5/3$.

Put $m=1-9d/8$. The sixth- and third-moment main terms are
\[
 H=m^{11}(1-111d/16),\qquad F=m^5(1-129d/32).
\]
The negative mean substitution is valid: for the sixth bound on
$d\le3/32$, $2m_0-11(31/32)d\ge2-417d/32>0$;
for the cubic on $d\le4/19$, use $2-231d/32>0$.
Convexity follows from
\[
 [(1-cd)^n(1-bd)]''=nc(1-cd)^{n-2}[(n-1)c+2b-(n+1)bcd],
\]
with $(n,b)=(11,111/16)$ or $(5,129/32)$, $c=9/8$.
For $d\le1/36$ the exponential criterion applies with
$D\ge(4a+1-\tau)d+(a^2-a)d^2$:
$D-3U>d[11/32-49d/4]\ge d/288$.

The remaining sixth-moment data are
\begin{center}\small
\begin{tabular}{@{}cccccc@{}}\toprule
$I$&$\bar m$&$M$&$V_p/d$&$B_{12}$&endpoints\\\midrule
$[1/36,1/16]$&$31/32$&$9/32$&$9/16$&$21$&$191/192,255/256$\\
$[1/16,5/64]$&$119/128$&$13/36$&$3/5$&$17$&$573/640,319/320$\\
$[5/64,3/32]$&$11/12$&$13/36$&$13/20$&$29/2$&$2397/2560,1271/1280$\\\bottomrule
\end{tabular}
\end{center}
Here the second row may enlarge its row-norm bound $13/40$ to $13/36$
for the remainder. The displayed endpoints bound $H+dB_{12}(V_p/d)$.
For the cubic ranges the table gives a rounded coefficient $k$ with
$B_6(V_p/d)\le k$; the final columns bound $F+kd$.
\begin{center}\small
\begin{tabular}{@{}cccccc@{}}\toprule
$I$&$\bar m$&$M$&$V_p/d$&$k$&endpoints\\\midrule
$[3/32,1/8]$&$17/19$&$7/16$&$59/80$&$5/2$&$19/32,9/16$\\
$[1/8,5/32]$&$13/15$&$51/100$&$17/20$&$11/4$&$19/32,371/640$\\
$[5/32,3/16]$&$33/40$&$23/40$&$19/20$&$8/3$&$17/30,7/12$\\
$[3/16,4/19]$&$15/19$&$5/8$&$25/24$&$125/48$&$439/768,269/456$\\\bottomrule
\end{tabular}
\end{center}
The complete finite $B_6$ bounds used here are $27/8,16/5,14/5,5/2$.
All cubic endpoints are less than $3/5$, so the cap converts them to the
sixth moment. The finite sums $B_r$, increasing $U,W$, and convexity
certify all closed intervals. Together with the two outer ranges this
proves $\E[p(qf_b)^6]\le1$ for every $b\ge6$.
